\documentclass[11pt,a4paper]{article}

% ---------- Packages ----------
\usepackage[utf8]{inputenc}
\usepackage[T1]{fontenc}
\usepackage{lmodern}
\usepackage{amsmath,amssymb}
\usepackage{enumitem}
\usepackage{geometry}
\usepackage{fancyhdr}
\usepackage[most]{tcolorbox}   % enhanced, breakable
\usepackage{hyperref}

% ---------- Geometry ----------
\geometry{margin=2.5cm}

% ---------- Header/Footer ----------
\setlength{\headheight}{14pt}  % fancyhdr warning verhindern
\pagestyle{fancy}
\fancyhf{}
\lhead{Einführung in die Elektrotechnik}
\rhead{Wintersemester 2025}
\cfoot{\thepage}

% ---------- Exercise Environment ----------
\newcounter{exercise}
\newenvironment{exercise}[1][]
  {\refstepcounter{exercise}%
   \begin{tcolorbox}[enhanced, breakable, width=\textwidth,
     colback=blue!5!white,colframe=blue!60!black,
     fonttitle=\bfseries\large,
     title=Aufgabe~\theexercise:~#1]}
  {\end{tcolorbox}}

% ---------- Solution Environment ----------
\newif\ifshowsolutions
\showsolutionstrue  % comment to hide solutions

\newenvironment{solution}
  {\ifshowsolutions
     \begin{tcolorbox}[enhanced, breakable, width=\textwidth,
       colback=green!5!white,colframe=green!50!black,
       fonttitle=\bfseries\large,
       title=Lösung]}
  {\end{tcolorbox}\fi}

% ---------- Document ----------
\begin{document}

% ----- Header Box -----
\begin{tcolorbox}[enhanced, breakable, colback=gray!10!white,
  colframe=gray!70!black, fonttitle=\bfseries\large,
  title={Aufgabenblatt \#1}: Grundlegende Konzepte, sharp corners, width=\textwidth]
\begin{center}
  Besprechung: 15.10.2025
\end{center}
\end{tcolorbox}
\vspace{1cm}

% ----- Aufgabe 1 -----
\begin{exercise}[Coulomb-Kraft]
Zwei Punktladungen sind $60\,\text{cm}$ voneinander entfernt und tragen eine Gesamtladung von $200\,\mu\text{C}$. Bestimmen Sie die Ladung der beiden Teilchen, wenn sie
\begin{enumerate}[label=\alph*)]
    \item sich mit einer Kraft von $80\,\text{N}$ abstoßen,
    \item sich mit einer Kraft von $80\,\text{N}$ anziehen.
\end{enumerate}

\ifshowsolutions
\begin{solution}
Die Ladungen seien $q_1$ und $q_2$ mit
\[
q_1 + q_2 = Q = 200 \cdot 10^{-6}\,\text{C}.
\]

Die Coulomb-Kraft ist
\begin{align*}
F &= k \frac{|q_1 q_2|}{r^2},\\
k &= 8.99 \cdot 10^9 \,\text{N\,m}^2/\text{C}^2,\\
r &= 0.60\,\text{m},\\
F &= 80\,\text{N}.
\end{align*}

Also gilt
\[
|q_1 q_2| = \frac{F r^2}{k} \approx 3.2 \cdot 10^{-9}\,\text{C}^2,
\]
mit
\[
q_1 q_2 = q_1 (q - q_1).
\]

\begin{enumerate}[label=\alph*)]
\item
\[
q_1 q_2 = +3.2 \cdot 10^{-9} \Rightarrow x^2 - (200 \cdot 10^{-6})x + 3.2 \cdot 10^{-9} = 0
\]
Mit
\[
 a x^2 + b x + c = 0 \Rightarrow x = \frac{-b \pm \sqrt{b^2 - 4ac}}{2a},
\]
\[
q_1 \approx 1.8 \cdot 10^{-4}\,\text{C},\quad
q_2 \approx 1.8 \cdot 10^{-5}\,\text{C}.
\]

\item
\[
q_1 q_2 = -3.2 \cdot 10^{-9} \Rightarrow x^2 - (200 \cdot 10^{-6})x - 3.2 \cdot 10^{-9} = 0
\]
\[
q_1 \approx -1.5 \cdot 10^{-5}\,\text{C},\quad
q_2 \approx 2.1 \cdot 10^{-4}\,\text{C}.
\]
\end{enumerate}
\end{solution}
\fi
\end{exercise}

% ----- Aufgabe 2 -----
\begin{exercise}[Ohmsches Gesetz]
Ein Widerstand von $15\, \Omega$ liegt an einer Spannung von $30\, \mathrm{V}$. Wie groß sind Stromstärke und Leitwert?  

\ifshowsolutions
\begin{solution}
Ohmsches Gesetz:  
\[
I = \frac{U}{R}
\]

Einsetzen der Werte:  
\[
I = \frac{30\, \mathrm{V}}{15\, \Omega} = 2\, \mathrm{A}
\]

Der Leitwert ist der Kehrwert des Widerstands:  
\[
G = \frac{1}{R} = \frac{1}{15\, \Omega} \approx 0.067\, \mathrm{S}
\]
\end{solution}
\fi
\end{exercise}

% ----- Aufgabe 3 -----
\begin{exercise}[Elektrischer Widerstand]
Gegeben seien zwei Kupferdrähte von gleicher Masse und kreisförmigem Querschnitt; Draht A ist doppelt so lang wie Draht B. Wie groß sind die Widerstände der beiden Drähte im Vergleich (unter Vernachlässigung von Temperatureffekten)? 
\begin{enumerate}[label=\alph*)]
    \item $R_A = 8 R_B$,
    \item $R_A = 4 R_B$
    \item $R_A = 2 R_B$,
    \item $R_A = R_B$.
\end{enumerate}

\ifshowsolutions
\begin{solution}
Der Widerstand eines Drahtes ist gegeben durch
\[
R = \rho \frac{l}{A},
\]
wobei $\rho$ der spezifische Widerstand, $l$ die Länge und $A$ die Querschnittsfläche ist.

Da beide Drähte die gleiche Masse und das gleiche Material haben, gilt
\[
m = \rho_m V = \rho_m (A \cdot l) = \text{const.}
\]
Also ist $A \cdot l = \text{const.}$, was bedeutet, dass $A \propto \frac{1}{l}$.

Für Draht A ist die Länge $l_A = 2 l_B$. Daher ist die Querschnittsfläche
\[
A_A = \frac{1}{2} A_B.
\]

Somit ergibt sich das Widerstandsverhältnis zu
\[
\frac{R_A}{R_B} 
= \frac{\rho \frac{l_A}{A_A}}{\rho \frac{l_B}{A_B}}
= \frac{l_A}{l_B} \cdot \frac{A_B}{A_A}
= \frac{2l_B}{l_B} \cdot \frac{A_B}{\frac{1}{2}A_B}
= 2 \cdot 2
= 4.
\]

Daher gilt
\[
R_A = 4 R_B.
\]
\end{solution}
\fi
\end{exercise}

% ----- Aufgabe 4 -----
\begin{exercise}[Spezifischer Widerstand]
Angenommen, wir haben einen Kupferdraht mit einem Widerstand von $2 \,\Omega$. Wir möchten den Kupferdraht durch einen Aluminiumdraht gleicher Länge und gleichen Widerstands ersetzen. Um welchen Faktor muss der Durchmesser des Drahtes vergrößert werden?

\ifshowsolutions
\begin{solution}
Für einen zylindrischen Draht gilt
\[
A = \frac{\pi d^2}{4},
\]
wobei \(d\) der Durchmesser ist.

Wir verlangen, dass der Aluminiumdraht den gleichen Widerstand \(R\) wie der Kupferdraht hat. Außerdem haben beide Drähte die gleiche Länge \(l\). Somit gilt
\[
\rho_{\text{Cu}} \frac{l}{A_{\text{Cu}}} = \rho_{\text{Al}} \frac{l}{A_{\text{Al}}}.
\]

Kürzen von \(l\) ergibt
\[
\frac{A_{\text{Al}}}{A_{\text{Cu}}} = \frac{\rho_{\text{Al}}}{\rho_{\text{Cu}}}.
\]

Da \(A \propto d^2\),
\[
\left(\frac{d_{\text{Al}}}{d_{\text{Cu}}}\right)^2 = \frac{\rho_{\text{Al}}}{\rho_{\text{Cu}}}.
\]

Somit gilt
\[
\frac{d_{\text{Al}}}{d_{\text{Cu}}} = \sqrt{\frac{\rho_{\text{Al}}}{\rho_{\text{Cu}}}}.
\]

Numerische Werte für den spezifischen Widerstand:
\[
\rho_{\text{Cu}} \approx 0.017 \,\frac{\Omega \cdot \text{mm}^2}{\text{m}},
\quad 
\rho_{\text{Al}} \approx 0.027 \,\frac{\Omega \cdot \text{mm}^2}{\text{m}}
\]

Also ist
\[
\frac{d_{\text{Al}}}{d_{\text{Cu}}} = \sqrt{\frac{0.027}{0.017}} \approx 1.26.
\]

\[
\boxed{\text{Der Durchmesser des Aluminiumdrahtes muss etwa 1,3-mal größer sein.}}
\]
\end{solution}
\fi
\end{exercise}

% ----- Aufgabe 5 -----
\begin{exercise}[Temperaturabhängigkeit des Widerstands]
Bei welcher Temperatur hat ein Kupferdraht einen um 10\% größeren Widerstand als bei 20\,$^\circ$C?

\ifshowsolutions
\begin{solution}
Die Temperaturabhängigkeit des Widerstands wird beschrieben durch
\[
R(T) = R_{0}\bigl(1 + \alpha (T - T_{0})\bigr),
\]
wobei $R_{0}$ der Widerstand bei der Referenztemperatur $T_{0}$ und $\alpha$ der Temperaturkoeffizient des Widerstands ist.

Wir suchen die Temperatur $T$, bei der
\[
R(T) = 1.1 \, R_{0}.
\]

Also
\[
1.1 R_{0} = R_{0}\bigl(1 + \alpha (T - 20^\circ\mathrm{C})\bigr).
\]

Dividieren durch $R_{0}$ ergibt
\[
1.1 = 1 + \alpha (T - 20).
\]

Somit gilt
\[
T - 20 = \frac{0.1}{\alpha}.
\]

Für Kupfer ist $\alpha \approx 0.0039 \,\text{K}^{-1}$, also
\[
T \approx 20 + \frac{0.1}{0.0039} \approx 45.6^\circ\mathrm{C}.
\]

Demnach hat der Kupferdraht bei etwa
\[
\boxed{46^\circ\mathrm{C}}
\]
einen um 10\% größeren Widerstand als bei 20\,$^\circ$C.
\end{solution}
\fi
\end{exercise}

% ----- Aufgabe 6 -----
\begin{exercise}[Energie]
Eine wiederaufladbare Taschenlampenbatterie kann etwa $90 \,\mathrm{mA}$ für ca. $12 \,\mathrm{h}$ liefern. Wieviel Ladung kann sie bei dieser Stromstärke abgeben? Wenn ihre Klemmenspannung $1.5 \,\mathrm{V}$ beträgt, wieviel Energie kann die Batterie liefern?

\ifshowsolutions
\begin{solution}
Die Batterie liefert einen Strom von
\[
I = 0.09 \,\mathrm{A},
\]
für eine Zeit von
\[
t = 12 \cdot 3600 \,\mathrm{s} = 43200 \,\mathrm{s}.
\]

\begin{enumerate}[label=\alph*)]
\item \[
Q = I \cdot t = 0.09 \cdot 43200 \,\mathrm{C} = 3888 \,\mathrm{C}
\]

Also ist die gesamte Ladung
\[
Q \approx 3.9 \,\mathrm{kC}.
\]

\item Die Energie ist gegeben durch
\[
W = Q \cdot V,
\]
wobei \(V = 1.5 \,\mathrm{V}\). Somit ist
\[
W = 3888 \cdot 1.5 \,\mathrm{J} = 5832 \,\mathrm{J}.
\]

Die gesamte Energie beträgt also
\[
W \approx 5.8 \,\mathrm{kJ}.
\]
\end{enumerate}
\end{solution}
\fi
\end{exercise}

% ----- Aufgabe 7 -----
\begin{exercise}[Leistung]
Gesucht wird der Widerstand einer Glühbirne mit Nennwert $60\,\text{W}$ bei $120\,\text{V}$.

\ifshowsolutions
\begin{solution}
Die elektrische Leistung hängt von Spannung und Widerstand ab gemäß
\[
P = \frac{V^2}{R}.
\]

Umstellen nach \(R\) ergibt
\[
R = \frac{V^2}{P}.
\]

Einsetzen der gegebenen Werte:
\[
V = 120\,\text{V}, 
\quad 
P = 60\,\text{W}
\]

\[
\Rightarrow R = \frac{(120)^2}{60} \,\Omega = \frac{14400}{60} \,\Omega = 240 \,\Omega
\]

Der \emph{heiße Widerstand} der Glühbirne beträgt also \(240 \,\Omega\).

\bigskip

\textbf{Bemerkung:}
Der Glühfaden einer Glühbirne hat im kalten Zustand einen viel geringeren Widerstand, typischerweise etwa
\[
R_{\text{cold}} \approx \frac{R_{\text{hot}}}{10} \sim 20\text{--}30 \,\Omega,
\]
abhängig vom Material und dem Temperaturkoeffizienten des Widerstands. Dies erklärt den hohen Strom, wenn die Birne erstmals eingeschaltet wird.
\end{solution}
\fi
\end{exercise}

% ----- Aufgabe 8 -----
\begin{exercise}[Teilchen]
Ein Teilchenbeschleuniger erzeugt einen Protonenstrahl mit einer Stromstärke von $3.5 \,\mu\mathrm{A}$. Jedes Proton hat eine Energie von $60 \,\mathrm{MeV}$. Die Protonen treffen auf ein Ziel in einer Vakuumkammer und kommen darin zur Ruhe, wodurch das Ziel erhitzt wird.
\begin{enumerate}[label=\alph*)]
    \item Wie viele Protonen treffen pro Sekunde auf das Ziel?
    \item Wieviel Energie wird pro Sekunde im Ziel umgesetzt?
\end{enumerate}

\ifshowsolutions
\begin{solution}
Die Stromstärke des Strahls beträgt $I = 3.5 \,\mu\mathrm{A} = 3.5 \cdot 10^{-6} \,\mathrm{A}$.  
Die Elementarladung ist $e = 1.602 \cdot 10^{-19} \,\mathrm{C}$.  

\begin{enumerate}[label=\alph*)]
\item Die Anzahl der Protonen pro Sekunde ist
\[
\dot{N}=\frac{I}{e}.
\]
Einsetzen ergibt
\[
\dot{N} = \frac{3.5 \cdot 10^{-6}}{1.602 \cdot 10^{-19}} \,\mathrm{s}^{-1} \approx 2.18 \cdot 10^{13} \,\mathrm{s}^{-1}.
\]
\item Jedes Proton hat eine Energie von $60 \,\mathrm{MeV}$.

Die pro Sekunde umgesetzte Energie beträgt
\[
P = V \cdot I = \frac{W}{Q} \cdot I = \frac{60 \,\mathrm{MeV}}{e} \cdot 3.5 \,\mu\mathrm{A} = 210 \,\mathrm{W}.
\]
\end{enumerate}
\end{solution}
\fi
\end{exercise}

\end{document}